\documentclass[11pt]{article}
\usepackage[a4paper,margin=0.82in]{geometry}
\usepackage{amsmath,amssymb,bm}
\usepackage{graphicx}
\usepackage{booktabs}
\usepackage{tabularx}
\usepackage{array}
\usepackage{microtype}
\usepackage{hyperref}
\usepackage{xcolor}
\usepackage{enumitem}
\usepackage{caption}
\usepackage{float}
\usepackage{fancyhdr}
\usepackage{titlesec}
\usepackage{setspace}
\definecolor{navy}{HTML}{17365D}
\hypersetup{colorlinks=true,linkcolor=navy,citecolor=navy,urlcolor=navy}
\titleformat{\section}{\Large\bfseries\color{navy}}{\thesection}{0.6em}{}
\titleformat{\subsection}{\large\bfseries}{\thesubsection}{0.5em}{}
\setlist[itemize]{leftmargin=1.35em,itemsep=2pt,topsep=3pt}
\setlength{\parindent}{0pt}
\setlength{\parskip}{5pt}
\captionsetup{font=small,labelfont=bf}
\setlength{\headheight}{14pt}
\pagestyle{fancy}
\fancyhf{}
\lhead{\CourseCode\ -- \ReportType}
\rhead{\StudentID}
\cfoot{\thepage}
\newcommand{\code}[1]{\texttt{#1}}


% ================= EDIT REPORT DETAILS HERE =================
% English text: compile with pdfLaTeX (or XeLaTeX).
% Chinese text: compile with XeLaTeX and uncomment the next line.
% \usepackage[UTF8,fontset=fandol]{ctex}
\newcommand{\University}{National University of Singapore}
\newcommand{\CourseCode}{COURSE0000}
\newcommand{\CourseName}{Course Name}
\newcommand{\ReportType}{Project Report}
\newcommand{\ReportTitle}{Your Report Title}
\newcommand{\StudentName}{Your Name}
\newcommand{\StudentID}{Your Student ID}
\newcommand{\StudentEmail}{your.email@example.com}
\newcommand{\ReportDate}{Month Year}
% ============================================================

\begin{document}
\begin{titlepage}
\centering
\vspace*{2.0cm}
{\Large \University\par}
\vspace{0.4cm}
{\large \CourseCode\ \CourseName\par}
\vspace{2.0cm}
{\Huge\bfseries \ReportType\par}
\vspace{0.35cm}
{\LARGE \ReportTitle\par}
\vfill
\begin{tabular}{rl}
\textbf{Name:} & \StudentName\\
\textbf{Student ID:} & \StudentID\\
\textbf{Email:} & \StudentEmail\\
\end{tabular}
\vfill
{\large \ReportDate\par}
\end{titlepage}

\begin{abstract}
[Summarize the problem, approach, main findings, and conclusion.
Replace this paragraph with your own abstract.]
\end{abstract}

% Optional contents page: uncomment both lines for longer reports.
% \tableofcontents
% \newpage

\section{Introduction}
[Introduce the topic, explain the motivation, and state the objectives
or questions addressed by this report.]

\section{Background and Notation}
[Summarize relevant concepts and prior work. Define the notation used
throughout the report. Cite sources where appropriate, e.g.,
\cite{example-reference}.]

\section{Methodology}
\subsection{Data or Problem Setup}
[Describe the data, assumptions, experimental setup, or problem context.]

\subsection{Approach}
[Explain your method and justify the key choices.]
An example numbered equation is
\begin{equation}
\widehat{y}_i = \bm{x}_i^{\mathsf T}\bm{\beta}.
\label{eq:example}
\end{equation}
Refer to it using Equation~\ref{eq:example}.

\section{Results and Discussion}
[Present the evidence, explain the findings, and discuss limitations.
The values below are examples only; replace or remove them.]

\begin{table}[htbp]
\centering
\caption{Example comparison table. Replace all example entries.}
\label{tab:example}
\small
\begin{tabularx}{\textwidth}{l r X}
\toprule
Method & Score & Notes \\
\midrule
Method A & 0.80 & Brief explanation of the first result. \\
Method B & 0.85 & Brief explanation of the second result. \\
\bottomrule
\end{tabularx}
\end{table}

Table~\ref{tab:example} demonstrates a table with an automatically
wrapping text column.

% Replace the placeholder with an image when ready. Paths are relative
% to this .tex file; create a figures/ folder beside it.
\begin{figure}[htbp]
\centering
% \includegraphics[width=0.88\textwidth]{figures/your-figure.png}
\fbox{\parbox[c][4cm][c]{0.84\textwidth}{\centering
Figure placeholder\\[0.5em]
Replace this box with your own figure.}}
\caption{Example figure caption. Describe the figure and its key message.}
\label{fig:example}
\end{figure}

Refer to figures using Figure~\ref{fig:example}. Use \code{filename.py}
for short code or file references; escape underscores as \code{my\_file.py}.

\section{Conclusion}
[Answer the report's main question, summarize the supported findings,
and identify any useful next steps.]

% Include only if required by your course. Describe your actual AI use
% accurately and follow the course's declaration requirements.
% \section*{AI Use Declaration}
% [State which tools you used, for which tasks, and how you checked their output.]

\begin{thebibliography}{9}
\bibitem{example-reference}
Author Name.
\newblock Title of the referenced work.
\newblock \emph{Journal or Publisher}, Year.
\end{thebibliography}

% Optional appendix: uncomment and replace the placeholder.
% \appendix
% \section{Additional Details}
% [Provide supporting derivations, additional results, or implementation details.]

\end{document}
